\documentclass[11pt]{article}

\usepackage[margin=1in]{geometry}
\usepackage{amsmath, amssymb}
\usepackage{listings}
\usepackage{xcolor}
\usepackage[T1]{fontenc}
\usepackage{lmodern}
\usepackage{enumitem}
\usepackage{parskip}
\usepackage{hyperref}

\hypersetup{
    colorlinks=true,
    linkcolor=blue,
    urlcolor=blue,
    citecolor=blue
}

\lstset{
    language=Python,
    basicstyle=\ttfamily\small,
    keywordstyle=\color{blue},
    commentstyle=\color{teal},
    stringstyle=\color{red!60!black},
    showstringspaces=false,
    columns=fullflexible,
    keepspaces=true,
    frame=single,
    breaklines=true,
    aboveskip=0.5em,
    belowskip=0.5em
}

\title{CS273P Machine Learning\\Final Exam Practice Problems}
\date{}
\author{}

\begin{document}

\maketitle

\begin{center}
\fbox{\parbox{0.9\textwidth}{%
\centering
\textbf{Overview.} This handout contains sample coding questions similar to those you may see on the final exam. Each problem includes a solution. Try solving each on your own before reading the solution.
}}
\end{center}

\vspace{1em}

\noindent\textbf{Instructions}
\begin{itemize}[leftmargin=1.5em]
    \item Write your solutions in Python.
    \item You may use \texttt{numpy} and \texttt{math}.
    \item Unless otherwise stated, aim for concise implementations.
    \item Most problems are designed to require only a few lines of code.
    \item You may assume inputs are valid (e.g., non-empty lists, compatible dimensions).
\end{itemize}

\vspace{1.5em}
\hrule
\vspace{1em}

\noindent\textbf{\large Part I: K-Nearest Neighbors}
\vspace{0.5em}

\section*{1. Euclidean Distance}
Write a function that computes the Euclidean distance between two 1D vectors.

\begin{lstlisting}
def euclidean_distance(x: list[float], y: list[float]) -> float:
    pass
\end{lstlisting}

\medskip
\noindent\textbf{Solution}
\begin{lstlisting}
import math

def euclidean_distance(x: list[float], y: list[float]) -> float:
    return math.sqrt(sum((a - b) ** 2 for a, b in zip(x, y)))
\end{lstlisting}

\section*{2. Nearest Neighbor}
Given a query vector and a list of training vectors, return the index of the nearest training example using Euclidean distance.

\begin{lstlisting}
def nearest_neighbor(X_train: list[list[float]], query: list[float]) -> int:
    pass
\end{lstlisting}

\medskip
\noindent\textbf{Solution}
\begin{lstlisting}
def nearest_neighbor(X_train: list[list[float]], query: list[float]) -> int:
    return min(range(len(X_train)), key=lambda i: euclidean_distance(X_train[i], query))
\end{lstlisting}

\section*{3. k Nearest Neighbor Indices}
Return the indices of the \texttt{k} nearest training examples to a query point, ordered from nearest to farthest.

\begin{lstlisting}
def knn_indices(X_train: list[list[float]], query: list[float], k: int) -> list[int]:
    pass
\end{lstlisting}

\medskip
\noindent\textbf{Solution}
\begin{lstlisting}
def knn_indices(X_train: list[list[float]], query: list[float], k: int) -> list[int]:
    return sorted(range(len(X_train)), key=lambda i: euclidean_distance(X_train[i], query))[:k]
\end{lstlisting}

\section*{4. KNN Majority Vote}
Given training data, binary labels, a query point, and \texttt{k}, predict the label (0 or 1) using majority vote among the \texttt{k} nearest neighbors.

\begin{lstlisting}
def knn_predict(X_train: list[list[float]], y_train: list[int], query: list[float], k: int) -> int:
    pass
\end{lstlisting}

\medskip
\noindent\textbf{Solution}
\begin{lstlisting}
def knn_predict(X_train: list[list[float]], y_train: list[int], query: list[float], k: int) -> int:
    idx = knn_indices(X_train, query, k)
    votes = sum(y_train[i] for i in idx)
    return 1 if votes >= k - votes else 0
\end{lstlisting}

\vspace{0.75em}
\noindent\textbf{\large Part II: Naive Bayes}
\vspace{0.5em}

\section*{5. Class Priors}
Given a list of class labels, compute the prior probability of each class and return a dictionary.

\begin{lstlisting}
def class_priors(y: list[int]) -> dict[int, float]:
    pass
\end{lstlisting}

\medskip
\noindent\textbf{Solution}
\begin{lstlisting}
def class_priors(y: list[int]) -> dict[int, float]:
    out = {}
    n = len(y)
    for c in y:
        out[c] = out.get(c, 0) + 1
    for c in out:
        out[c] /= n
    return out
\end{lstlisting}

\section*{6. Bernoulli Naive Bayes Feature Probability}
Assume binary features and class labels. Compute
\[
P(x_j = 1 \mid y = c)
\]
for feature index \texttt{j} and class \texttt{c}. Assume at least one training example has class \texttt{c}.

\begin{lstlisting}
def bernoulli_feature_prob(X: list[list[int]], y: list[int], j: int, c: int) -> float:
    pass
\end{lstlisting}

\medskip
\noindent\textbf{Solution}
\begin{lstlisting}
def bernoulli_feature_prob(X: list[list[int]], y: list[int], j: int, c: int) -> float:
    vals = [X[i][j] for i in range(len(y)) if y[i] == c]
    return sum(vals) / len(vals)
\end{lstlisting}

\section*{7. Naive Bayes Log Score}
Given a binary input vector \texttt{x}, class prior \texttt{p\_y}, and Bernoulli probabilities \texttt{p\_x\_given\_y}, compute the Naive Bayes log score up to a constant.

\begin{lstlisting}
def nb_log_score(x: list[int], p_y: float, p_x_given_y: list[float]) -> float:
    pass
\end{lstlisting}

\medskip
\noindent\textbf{Solution}
\begin{lstlisting}
import math

def nb_log_score(x: list[int], p_y: float, p_x_given_y: list[float]) -> float:
    s = math.log(p_y)
    for xj, pj in zip(x, p_x_given_y):
        s += math.log(pj) if xj == 1 else math.log(1 - pj)
    return s
\end{lstlisting}

\section*{8. Naive Bayes Prediction from Scores}
Given the log scores for two classes, return the predicted class.

\begin{lstlisting}
def predict_from_scores(score0: float, score1: float) -> int:
    pass
\end{lstlisting}

\medskip
\noindent\textbf{Solution}
\begin{lstlisting}
def predict_from_scores(score0: float, score1: float) -> int:
    return 1 if score1 > score0 else 0
\end{lstlisting}

\vspace{0.75em}
\noindent\textbf{\large Part III: Linear Regression}
\vspace{0.5em}

\section*{9. Mean Squared Error}
Compute the mean squared error between \texttt{y\_true} and \texttt{y\_pred}.

\begin{lstlisting}
def mse_loss(y_true: list[float], y_pred: list[float]) -> float:
    pass
\end{lstlisting}

\medskip
\noindent\textbf{Solution}
\begin{lstlisting}
def mse_loss(y_true: list[float], y_pred: list[float]) -> float:
    n = len(y_true)
    return sum((a - b) ** 2 for a, b in zip(y_true, y_pred)) / n
\end{lstlisting}

\section*{10. Linear Regression Prediction}
Given \texttt{x}, weight vector \texttt{w}, and bias \texttt{b}, compute the linear regression prediction
\[
\hat{y} = w^T x + b
\]

\begin{lstlisting}
def linear_predict(x: list[float], w: list[float], b: float) -> float:
    pass
\end{lstlisting}

\medskip
\noindent\textbf{Solution}
\begin{lstlisting}
def linear_predict(x: list[float], w: list[float], b: float) -> float:
    return sum(xj * wj for xj, wj in zip(x, w)) + b
\end{lstlisting}

\section*{11. One Gradient Descent Step for Linear Regression}
For one data point \texttt{(x, y)}, perform one gradient descent update for \texttt{w} and \texttt{b} using squared loss.

\begin{lstlisting}
def linear_regression_step(
    x: list[float],
    y: float,
    w: list[float],
    b: float,
    lr: float
) -> tuple[list[float], float]:
    pass
\end{lstlisting}

\medskip
\noindent\textbf{Solution}
\begin{lstlisting}
def linear_regression_step(
    x: list[float],
    y: float,
    w: list[float],
    b: float,
    lr: float
) -> tuple[list[float], float]:
    y_hat = linear_predict(x, w, b)
    err = y_hat - y
    new_w = [wj - lr * 2 * err * xj for wj, xj in zip(w, x)]
    new_b = b - lr * 2 * err
    return new_w, new_b
\end{lstlisting}

\section*{12. Batch Gradient for Linear Regression}
Given a dataset \texttt{(X, y)}, compute the gradient of mean squared error with respect to \texttt{w} and \texttt{b}.

\begin{lstlisting}
def linear_regression_gradient(
    X: list[list[float]],
    y: list[float],
    w: list[float],
    b: float
) -> tuple[list[float], float]:
    pass
\end{lstlisting}

\medskip
\noindent\textbf{Solution}
\begin{lstlisting}
def linear_regression_gradient(
    X: list[list[float]],
    y: list[float],
    w: list[float],
    b: float
) -> tuple[list[float], float]:
    n = len(X)
    d = len(w)
    grad_w = [0.0] * d
    grad_b = 0.0

    for xi, yi in zip(X, y):
        err = linear_predict(xi, w, b) - yi
        for j in range(d):
            grad_w[j] += 2 * err * xi[j] / n
        grad_b += 2 * err / n

    return grad_w, grad_b
\end{lstlisting}

\vspace{0.75em}
\noindent\textbf{\large Part IV: Logistic Regression}
\vspace{0.5em}

\section*{13. Sigmoid}
Implement the sigmoid function
\[
\sigma(z) = \frac{1}{1 + e^{-z}}
\]

\begin{lstlisting}
def sigmoid(z: float) -> float:
    pass
\end{lstlisting}

\medskip
\noindent\textbf{Solution}
\begin{lstlisting}
import math

def sigmoid(z: float) -> float:
    return 1 / (1 + math.exp(-z))
\end{lstlisting}

\section*{14. Binary Negative Log Likelihood}
For binary label \texttt{y} in \(\{0, 1\}\) and predicted probability \texttt{p} (in \((0, 1)\)), compute the negative log likelihood (cross-entropy)
\[
-\bigl(y \log p + (1-y)\log(1-p)\bigr).
\]

\begin{lstlisting}
def binary_nll(y: int, p: float) -> float:
    pass
\end{lstlisting}

\medskip
\noindent\textbf{Solution}
\begin{lstlisting}
import math

def binary_nll(y: int, p: float) -> float:
    return -(y * math.log(p) + (1 - y) * math.log(1 - p))
\end{lstlisting}

\section*{15. Logistic Regression Probability}
Given \texttt{x}, \texttt{w}, and \texttt{b}, compute
\[
P(y=1 \mid x) = \sigma(w^T x + b)
\]

\begin{lstlisting}
def logistic_predict_prob(x: list[float], w: list[float], b: float) -> float:
    pass
\end{lstlisting}

\medskip
\noindent\textbf{Solution}
\begin{lstlisting}
def logistic_predict_prob(x: list[float], w: list[float], b: float) -> float:
    z = sum(xj * wj for xj, wj in zip(x, w)) + b
    return sigmoid(z)
\end{lstlisting}

\section*{16. Logistic Regression Gradient}
For one data point \texttt{(x, y)}, compute the gradient of binary negative log likelihood with respect to \texttt{w} and \texttt{b}.

\begin{lstlisting}
def logistic_gradient(
    x: list[float],
    y: int,
    w: list[float],
    b: float
) -> tuple[list[float], float]:
    pass
\end{lstlisting}

\medskip
\noindent\textbf{Solution}
\begin{lstlisting}
def logistic_gradient(
    x: list[float],
    y: int,
    w: list[float],
    b: float
) -> tuple[list[float], float]:
    p = logistic_predict_prob(x, w, b)
    grad_w = [(p - y) * xj for xj in x]
    grad_b = p - y
    return grad_w, grad_b
\end{lstlisting}

\vspace{0.75em}
\noindent\textbf{\large Part V: Multiclass Classification}
\vspace{0.5em}

\section*{17. Softmax}
Given a list of logits, return the softmax probabilities.

\begin{lstlisting}
def softmax(logits: list[float]) -> list[float]:
    pass
\end{lstlisting}

\medskip
\noindent\textbf{Solution}
\begin{lstlisting}
import math

def softmax(logits: list[float]) -> list[float]:
    exps = [math.exp(z) for z in logits]
    s = sum(exps)
    return [e / s for e in exps]
\end{lstlisting}

\section*{18. Multiclass Negative Log Likelihood}
Given a vector of class probabilities and the correct class index \texttt{y}, compute
\[
-\log p_y
\]

\begin{lstlisting}
def multiclass_nll(probs: list[float], y: int) -> float:
    pass
\end{lstlisting}

\medskip
\noindent\textbf{Solution}
\begin{lstlisting}
import math

def multiclass_nll(probs: list[float], y: int) -> float:
    return -math.log(probs[y])
\end{lstlisting}

\section*{19. Predict Class from Logits}
Given a list of logits, return the predicted class index.

\begin{lstlisting}
def predict_class(logits: list[float]) -> int:
    pass
\end{lstlisting}

\medskip
\noindent\textbf{Solution}
\begin{lstlisting}
def predict_class(logits: list[float]) -> int:
    return max(range(len(logits)), key=lambda i: logits[i])
\end{lstlisting}

\vspace{0.75em}
\noindent\textbf{\large Part VI: Regularization}
\vspace{0.5em}

\section*{20. L2 Regularization Penalty}
Given a weight vector \texttt{w}, compute the L2 penalty
\[
\sum_j w_j^2
\]

\begin{lstlisting}
def l2_penalty(w: list[float]) -> float:
    pass
\end{lstlisting}

\medskip
\noindent\textbf{Solution}
\begin{lstlisting}
def l2_penalty(w: list[float]) -> float:
    return sum(wj ** 2 for wj in w)
\end{lstlisting}

\section*{21. L1 Regularization Penalty}
Given a weight vector \texttt{w}, compute the L1 penalty
\[
\sum_j |w_j|
\]

\begin{lstlisting}
def l1_penalty(w: list[float]) -> float:
    pass
\end{lstlisting}

\medskip
\noindent\textbf{Solution}
\begin{lstlisting}
def l1_penalty(w: list[float]) -> float:
    return sum(abs(wj) for wj in w)
\end{lstlisting}

\vspace{0.75em}
\noindent\textbf{\large Part VII: Neural Networks}
\vspace{0.5em}

\section*{22. ReLU Activation}
Apply ReLU elementwise to a vector.

\begin{lstlisting}
def relu(x: list[float]) -> list[float]:
    pass
\end{lstlisting}

\medskip
\noindent\textbf{Solution}
\begin{lstlisting}
def relu(x: list[float]) -> list[float]:
    return [max(0, v) for v in x]
\end{lstlisting}

\section*{23. Dense Layer Forward Pass}
Given input vector \texttt{x}, weight matrix \texttt{W}, and bias vector \texttt{b}, compute the output of a fully connected layer:
\[
Wx + b
\]
Assume \texttt{W[i]} is the weight row for output neuron \texttt{i}.

\begin{lstlisting}
def dense_forward(x: list[float], W: list[list[float]], b: list[float]) -> list[float]:
    pass
\end{lstlisting}

\medskip
\noindent\textbf{Solution}
\begin{lstlisting}
def dense_forward(x: list[float], W: list[list[float]], b: list[float]) -> list[float]:
    return [sum(wij * xj for wij, xj in zip(row, x)) + bi for row, bi in zip(W, b)]
\end{lstlisting}

\section*{24. 2D Valid Convolution}
Given a 2D image and a 2D kernel, compute the valid convolution output with stride 1 and no padding.

\begin{lstlisting}
def conv2d_valid(image: list[list[float]], kernel: list[list[float]]) -> list[list[float]]:
    pass
\end{lstlisting}

\medskip
\noindent\textbf{Solution}
\begin{lstlisting}
def conv2d_valid(image: list[list[float]], kernel: list[list[float]]) -> list[list[float]]:
    h, w = len(image), len(image[0])
    kh, kw = len(kernel), len(kernel[0])
    out = []

    for i in range(h - kh + 1):
        row = []
        for j in range(w - kw + 1):
            s = 0
            for a in range(kh):
                for b in range(kw):
                    s += image[i + a][j + b] * kernel[a][b]
            row.append(s)
        out.append(row)

    return out
\end{lstlisting}

\section*{25. 2x2 Max Pooling}
Given a 2D feature map, perform 2x2 max pooling with stride 2. Assume height and width are even.

\begin{lstlisting}
def max_pool_2x2(feature_map: list[list[float]]) -> list[list[float]]:
    pass
\end{lstlisting}

\medskip
\noindent\textbf{Solution}
\begin{lstlisting}
def max_pool_2x2(feature_map: list[list[float]]) -> list[list[float]]:
    h, w = len(feature_map), len(feature_map[0])
    out = []

    for i in range(0, h, 2):
        row = []
        for j in range(0, w, 2):
            block = [
                feature_map[i][j],
                feature_map[i][j + 1],
                feature_map[i + 1][j],
                feature_map[i + 1][j + 1],
            ]
            row.append(max(block))
        out.append(row)

    return out
\end{lstlisting}

\vspace{0.75em}
\noindent\textbf{\large Part VIII: RNNs, SVMs, and Transformers}
\vspace{0.5em}

\section*{26. One RNN Hidden-State Update}
Given input vector \texttt{x}, previous hidden state \texttt{h\_prev}, weight matrices \texttt{Wx}, \texttt{Wh}, and bias \texttt{b}, compute one RNN hidden-state update using
\[
h = \tanh(W_x x + W_h h_{prev} + b)
\]

\begin{lstlisting}
def rnn_step(
    x: list[float],
    h_prev: list[float],
    Wx: list[list[float]],
    Wh: list[list[float]],
    b: list[float]
) -> list[float]:
    pass
\end{lstlisting}

\medskip
\noindent\textbf{Solution}
\begin{lstlisting}
import math

def rnn_step(
    x: list[float],
    h_prev: list[float],
    Wx: list[list[float]],
    Wh: list[list[float]],
    b: list[float]
) -> list[float]:
    out = []
    for i in range(len(b)):
        s = b[i]
        s += sum(Wx[i][j] * x[j] for j in range(len(x)))
        s += sum(Wh[i][j] * h_prev[j] for j in range(len(h_prev)))
        out.append(math.tanh(s))
    return out
\end{lstlisting}

\section*{27. Hinge Loss}
For binary classification with label \texttt{y} in \(\{-1, +1\}\) and model score \texttt{s}, compute the hinge loss
\[
\max(0, 1 - y s)
\]

\begin{lstlisting}
def hinge_loss(y: int, s: float) -> float:
    pass
\end{lstlisting}

\medskip
\noindent\textbf{Solution}
\begin{lstlisting}
def hinge_loss(y: int, s: float) -> float:
    return max(0.0, 1 - y * s)
\end{lstlisting}

\section*{28. Scaled Dot-Product Attention Score}
Given query vector \texttt{q} and key vector \texttt{k}, compute the scaled dot-product attention score
\[
\frac{q \cdot k}{\sqrt{d}}
\]
where \texttt{d} is the vector dimension.

\begin{lstlisting}
def attention_score(q: list[float], k: list[float]) -> float:
    pass
\end{lstlisting}

\medskip
\noindent\textbf{Solution}
\begin{lstlisting}
import math

def attention_score(q: list[float], k: list[float]) -> float:
    d = len(q)
    return sum(a * b for a, b in zip(q, k)) / math.sqrt(d)
\end{lstlisting}

\section*{29. Scaled Dot-Product Attention Matrix}
Given a list of query vectors \texttt{Q} and a list of key vectors \texttt{K} (each vector of dimension \texttt{d}), compute the attention matrix of shape \texttt{(n\_queries, n\_keys)}. For each query, compute scaled dot-product scores with all keys, then apply softmax over the keys so each row sums to 1:
\[
A_{ij} = \frac{\exp(q_i \cdot k_j / \sqrt{d})}{\sum_{j'} \exp(q_i \cdot k_{j'} / \sqrt{d})}
\]

\begin{lstlisting}
def attention_matrix(Q: list[list[float]], K: list[list[float]]) -> list[list[float]]:
    pass
\end{lstlisting}

\medskip
\noindent\textbf{Solution}
\begin{lstlisting}
import math

def attention_matrix(Q: list[list[float]], K: list[list[float]]) -> list[list[float]]:
    d = len(Q[0])
    scale = math.sqrt(d)
    out = []
    for q in Q:
        scores = [sum(qi * kj for qi, kj in zip(q, k)) / scale for k in K]
        exps = [math.exp(s) for s in scores]
        total = sum(exps)
        out.append([e / total for e in exps])
    return out
\end{lstlisting}

\section*{30. Next Token Prediction}
Given a vector of logits over the vocabulary, return the index of the most likely next token.

\begin{lstlisting}
def next_token(logits: list[float]) -> int:
    pass
\end{lstlisting}

\medskip
\noindent\textbf{Solution}
\begin{lstlisting}
def next_token(logits: list[float]) -> int:
    return max(range(len(logits)), key=lambda i: logits[i])
\end{lstlisting}

\vspace{0.75em}
\noindent\textbf{\large Part IX: Clustering and PCA}
\vspace{0.5em}

\section*{31. Clustering / PCA Short Questions}

\subsection*{31(a) K-Means Assignment}
Given a point and a list of cluster centers, return the index of the nearest center.

\begin{lstlisting}
def assign_cluster(point: list[float], centers: list[list[float]]) -> int:
    pass
\end{lstlisting}

\medskip
\noindent\textbf{Solution}
\begin{lstlisting}
def assign_cluster(point: list[float], centers: list[list[float]]) -> int:
    return min(range(len(centers)), key=lambda i: euclidean_distance(point, centers[i]))
\end{lstlisting}

\subsection*{31(b) K-Means Center Update}
Given all points assigned to one cluster, return the new center as the mean point.

\begin{lstlisting}
def update_center(points: list[list[float]]) -> list[float]:
    pass
\end{lstlisting}

\medskip
\noindent\textbf{Solution}
\begin{lstlisting}
def update_center(points: list[list[float]]) -> list[float]:
    n = len(points)
    d = len(points[0])
    return [sum(points[i][j] for i in range(n)) / n for j in range(d)]
\end{lstlisting}

\subsection*{31(c) Gaussian Mixture Model: E-step and M-step}

\textbf{E-step.} Given mixture weights \texttt{pi1}, \texttt{pi2} and the likelihoods \texttt{p1}, \texttt{p2} of one data point under each component, compute the responsibilities (soft assignments) for that point:
\[
\gamma_1 = \frac{\pi_1 p_1}{\pi_1 p_1 + \pi_2 p_2}, \qquad \gamma_2 = 1 - \gamma_1.
\]
Return \texttt{($\gamma_1$, $\gamma_2$)}.

\begin{lstlisting}
def gmm_e_step(pi1: float, p1: float, pi2: float, p2: float) -> tuple[float, float]:
    pass
\end{lstlisting}

\medskip
\noindent\textbf{Solution}
\begin{lstlisting}
def gmm_e_step(pi1: float, p1: float, pi2: float, p2: float) -> tuple[float, float]:
    denom = pi1 * p1 + pi2 * p2
    return (pi1 * p1 / denom, pi2 * p2 / denom)
\end{lstlisting}

\textbf{M-step.} Given 1D data \texttt{X} and a list of responsibilities \texttt{resp} (each \texttt{resp[i]} is \texttt{($\gamma_{i1}$, $\gamma_{i2}$)} for point \texttt{X[i]}), compute the new mixture weights and component means:
\[
\pi_k = \frac{1}{n}\sum_{i=1}^n \gamma_{ik}, \qquad \mu_k = \frac{\sum_{i=1}^n \gamma_{ik} x_i}{\sum_{i=1}^n \gamma_{ik}}.
\]
Return \texttt{(pi1\_new, pi2\_new, mu1\_new, mu2\_new)}.

\begin{lstlisting}
def gmm_m_step(X: list[float], resp: list[tuple[float, float]]) -> tuple[float, float, float, float]:
    pass
\end{lstlisting}

\medskip
\noindent\textbf{Solution}
\begin{lstlisting}
def gmm_m_step(X: list[float], resp: list[tuple[float, float]]) -> tuple[float, float, float, float]:
    n = len(X)
    g1 = [r[0] for r in resp]
    g2 = [r[1] for r in resp]
    n1, n2 = sum(g1), sum(g2)
    pi1_new = n1 / n
    pi2_new = n2 / n
    mu1_new = sum(g1[i] * X[i] for i in range(n)) / n1 if n1 > 0 else 0.0
    mu2_new = sum(g2[i] * X[i] for i in range(n)) / n2 if n2 > 0 else 0.0
    return (pi1_new, pi2_new, mu1_new, mu2_new)
\end{lstlisting}

\subsection*{31(d) Mean Centering for PCA}
Given a data matrix \texttt{X}, subtract the mean of each feature column.

\begin{lstlisting}
def center_data(X: list[list[float]]) -> list[list[float]]:
    pass
\end{lstlisting}

\medskip
\noindent\textbf{Solution}
\begin{lstlisting}
def center_data(X: list[list[float]]) -> list[list[float]]:
    n = len(X)
    d = len(X[0])
    means = [sum(X[i][j] for i in range(n)) / n for j in range(d)]
    return [[X[i][j] - means[j] for j in range(d)] for i in range(n)]
\end{lstlisting}

\subsection*{31(e) PCA via Eigenvalue Decomposition in NumPy}
Given a mean-centered data matrix \texttt{X} as a NumPy array of shape \texttt{(n, d)}, compute the first principal component, that is, the eigenvector corresponding to the largest eigenvalue of the covariance matrix.

\begin{lstlisting}
def first_principal_component(X):
    pass
\end{lstlisting}

\medskip
\noindent\textbf{Solution}
\begin{lstlisting}
import numpy as np

def first_principal_component(X):
    C = np.cov(X, rowvar=False)
    eigvals, eigvecs = np.linalg.eig(C)
    idx = np.argmax(np.real(eigvals))
    return np.real(eigvecs[:, idx])
\end{lstlisting}

\end{document}
